\begin{lemma}
For every \(\varepsilon>0\), put
\[
M_n=\operatorname{TwoBiteNaturalReducedVertexCount}(n),\qquad
p_n=\operatorname{TwoBiteEdgeProbability}(1/2,n),\qquad
c_\varepsilon=\varepsilon^2/100 .
\]
Then
\[
M_n\bigl(\exp(-c_\varepsilon p_nM_n)+
          \exp(-c_\varepsilon p_nM_n)\bigr)\to 0
\]
as \(n\to\infty\).
\end{lemma}
\begin{proof}
The definitions of \(p_n\) and \(M_n\) are
\noderef{TwoBiteEdgeProbability} and
\noderef{TwoBiteNaturalReducedVertexCount}.  By
\noderef{NaturalDegreeMassDominatesLogCube}, applied with
\(A=c_\varepsilon^{-1}\), eventually
\[
(\log n)^3\leq c_\varepsilon p_nM_n .
\]
Therefore each exponential term is at most \(\exp(-(\log n)^3)\).  Also
\(M_n\leq n\) eventually, since \(M_n\) is the floor of \(n/(\log n)^2\);
this uses \noderef{TwoBiteReducedVertexCount} and \noderef{TwoBiteStretch}.
Thus eventually
\[
M_n\bigl(\exp(-c_\varepsilon p_nM_n)+
          \exp(-c_\varepsilon p_nM_n)\bigr)
 \leq
2n\exp(-(\log n)^3)
=
\exp(\log 2+\log n-(\log n)^3).
\]
The exponent \(\log 2+y-y^3\) tends to \(-\infty\) as \(y\to\infty\), and
\(y=\log n\to\infty\).  Hence the right-hand side tends to \(0\), and the
squeeze theorem gives the claimed convergence.
\end{proof}
